\chapter{Classic Series and Infinite Products}

\begin{center}\itshape
$\dfrac{\pi^2}6=1+\dfrac14+\dfrac19+\dfrac1{16}+\cdots$\\[0.6em]
$\dfrac2\pi=\dfrac{\sqrt2}2\cdot\dfrac{\sqrt{2+\sqrt2}}2\cdot\dfrac{\sqrt{2+\sqrt{2+\sqrt2}}}2\cdots$
\end{center}

In 1735 the young Euler announced that $1+\frac14+\frac19+\cdots=\frac{\pi^2}6$, solving the Basel problem that had defeated
the Bernoullis. His argument was daring: the function $\frac{\sin x}x$ has the roots $\pm\pi,\pm2\pi,\dots$, so, like a
polynomial, it should be the product $\prod_n\bigl(1-\frac{x^2}{n^2\pi^2}\bigr)$; comparing the coefficients of $x^2$ on
both sides gives $-\frac16=-\sum_n\frac1{n^2\pi^2}$. Euler knew the step from polynomials to infinite products needed
justification, and spent years checking it by other routes.

This chapter does it properly. The key is a formula that Euler also found: the expansion of $\pi\cot\pi x$ as a sum
over its poles, proved here by Herglotz's trick, a little gem of an argument that uses only continuity. Integrating it
gives the product for the sine; expanding it in powers of $x$ and comparing with the Bernoulli numbers of Chapter~1
gives every $\zeta(2k)$ at once. Along the way come Wallis's and Vi\`ete's products for $\pi$, and the chapter ends with
the irrationality of $\pi$ and $\pi^2$.

Volume~1 proved $\zeta(2)=\frac{\pi^2}6$ with Vi\`ete's relations between the roots of a polynomial and promised five more
proofs; this chapter gives two, Chapter~8 two more with Fourier series, and Chapter~11 one with residues.

\begin{think}
The sine vanishes at every integer multiple of $\pi$. A polynomial with roots $a_1,\dots,a_n$ and value $1$ at $0$ is
$\prod\bigl(1-\frac x{a_i}\bigr)$. Is $\frac{\sin x}x$ determined by its roots? What about $e^x\frac{\sin x}x$?
\end{think}

\section{Infinite products}

\begin{definition}
Let $(u_n)_{n\ge1}$ be real numbers. The \emph{infinite product} $\prod_{n\ge1}u_n$ \emph{converges} if the partial products
$P_N=u_1\cdots u_N$ have a limit $P$, and its value is then $P$.
\end{definition}

A product with a zero factor converges to $0$; and $\prod\frac12$ converges to $0$ in this sense. The useful products
are those whose factors are close to $1$ and which converge to a non-zero limit.

\begin{proposition}\label{prop:products}
Let $a_n>-1$ be real numbers with $\sum|a_n|<\infty$. Then $\prod(1+a_n)$ converges to a limit in $\left]0,+\infty\right[$, and
$\prod(1+a_n)=\exp\sum\ln(1+a_n)$. If the $a_n=a_n(x)$ are continuous functions of $x$ in an interval and
$\sum_n\sup_x|a_n(x)|<\infty$, the product is a continuous function of $x$.
\end{proposition}
\begin{proof}
For $|a|\le\frac12$,
\[
\begin{aligned}|\ln(1+a)|&=\Bigl|\sum_{k\ge1}\frac{(-1)^{k-1}a^k}k\Bigr|\eqbr\le\sum_k|a|^k=\frac{|a|}{1-|a|}\le2|a|.\end{aligned}
\]
Since $a_n\to0$,
$|a_n|\le\frac12$ for $n\ge n_0$, so $\sum\ln(1+a_n)$ converges absolutely, and $P_N=\exp\sum_{n\le N}\ln(1+a_n)$ converges to
$\exp\sum_n\ln(1+a_n)>0$ by continuity of the exponential. In the second case the same bound and the Weierstrass test show that
$\sum\ln(1+a_n(x))$ converges uniformly in $x$ (for $n\ge n_0$, where $\sup|a_n|\le\frac12$), so its sum is continuous by
Lemma~\ref{lem:uniform}, and so is its exponential.
\end{proof}

\begin{example}
$\prod_{n\ge2}\bigl(1-\frac1{n^2}\bigr)=\frac12$, since $\prod_{n=2}^N\frac{(n-1)(n+1)}{n\cdot n}=\frac{N+1}{2N}$ telescopes.
\end{example}

\section{The partial fractions of the cotangent}

The function $\pi\cot\pi x$ has a simple pole with residue $1$ at every integer: near $n$, $\pi\cot\pi x\approx\frac1{x-n}$.
Euler guessed that it is the sum of these poles.

\begin{theorem}[Euler]\label{thm:cot}
For every real $x\notin\Z$,
\[
\begin{aligned}\pi\cot\pi x&=\frac1x+\sum_{n\ge1}\Bigl(\frac1{x+n}+\frac1{x-n}\Bigr)\eqbr=\frac1x+\sum_{n\ge1}\frac{2x}{x^2-n^2}.\end{aligned}
\]
\end{theorem}
\begin{proof}[Proof (Herglotz's trick)]
Let $f(x)=\pi\cot\pi x$ and $g(x)=\frac1x+\sum_{n\ge1}\frac{2x}{x^2-n^2}$, both defined on $\R\setminus\Z$. On a closed
bounded interval avoiding $\Z$, $\bigl|\frac{2x}{x^2-n^2}\bigr|\le\frac{2A}{n^2-A^2}$ for $n>A=\max|x|$, so the series converges
uniformly (Lemma~\ref{lem:uniform}) and $g$ is continuous. Write $g_N(x)=\sum_{n=-N}^N\frac1{x+n}$, so that $g=\lim g_N$. We compare four
properties of $f$ and $g$.

\emph{Both are odd.} Clear.

\emph{Both have period $1$.} For $f$ this is the periodicity of $\cot$. For $g$, $g_N(x+1)=\sum_{m=-N+1}^{N+1}\frac1{x+m}=g_N(x)+\frac1{x+N+1}-\frac1{x-N}$,
and the last two terms tend to $0$.

\emph{Both satisfy $\varphi\bigl(\frac x2\bigr)+\varphi\bigl(\frac{x+1}2\bigr)=2\varphi(x)$.} For $f$, put $y=\frac{\pi x}2$: since
$\cot\bigl(y+\frac\pi2\bigr)=-\tan y$,
\[
\begin{aligned}\cot y-\tan y&=\frac{\cos^2y-\sin^2y}{\sin y\cos y}\eqbr=\frac{2\cos2y}{\sin2y}=2\cot2y.\end{aligned}
\]
For $g$, let $G_N(x)=g_N\bigl(\frac x2\bigr)+g_N\bigl(\frac{x+1}2\bigr)$. Then
\[
\begin{aligned}G_N(x)&=\sum_{n=-N}^N\Bigl(\frac2{x+2n}+\frac2{x+2n+1}\Bigr)\eqbr=2g_{2N}(x)+\frac2{x+2N+1},\end{aligned}
\]
and letting $N\to\infty$ gives the relation.

\emph{Both are $\frac1x+o(1)$ at $0$.} For $g$, $g(x)-\frac1x=\sum_{n\ge1}\frac{2x}{x^2-n^2}\to0$ by uniform convergence near $0$.
For $f$, by Taylor's formula $\sin\pi x=\pi x-\frac{(\pi x)^3}6+O(x^5)$ and $\cos\pi x=1-\frac{(\pi x)^2}2+O(x^4)$, so
\[
\begin{aligned}\pi\cot\pi x-\frac1x&=\frac{\pi x\cos\pi x-\sin\pi x}{x\sin\pi x}\eqbr=\frac{-\frac{(\pi x)^3}3+O(x^5)}{\pi x^2+O(x^4)}\to0.\end{aligned}
\]

Now let $h=f-g$ on $\R\setminus\Z$. It is continuous, odd, of period $1$, satisfies the relation, and tends to $0$ at $0$,
hence at every integer by periodicity. Extend it by $h(n)=0$ for $n\in\Z$: it becomes continuous on $\R$, and the relation
$h(\frac x2)+h(\frac{x+1}2)=2h(x)$ extends to all $x$ by continuity. Being continuous and periodic, $h$ attains its maximum
$M$ at some $x_0$. Then $h(\frac{x_0}2)+h(\frac{x_0+1}2)=2M$ with both terms at most $M$, so $h(\frac{x_0}2)=M$. Repeating,
$h(\frac{x_0}{2^k})=M$ for every $k$, and letting $k\to\infty$, $M=h(0)=0$. The same argument applied to $-h$ shows that
the minimum is $0$. So $h=0$.
\end{proof}

The argument is a maximum principle in miniature: an averaging identity forces a continuous function to be constant.
Chapter~11 will reprove the theorem with residues, and Chapter~10 with Poisson summation.

\section{The product for the sine}

\begin{theorem}[Euler]\label{thm:sineproduct}
For every real $x$,
\[
\sin\pi x=\pi x\prod_{n\ge1}\Bigl(1-\frac{x^2}{n^2}\Bigr).
\]
\end{theorem}
\begin{proof}
By Proposition~\ref{prop:products}, the product $P(x)=\pi x\prod_{n\ge1}\bigl(1-\frac{x^2}{n^2}\bigr)$ converges for every $x$ (a
factor may vanish, at an integer; then $P(x)=0$), and it is continuous on $\left]-1,1\right[$, where all factors are positive.

Let $0<x<1$. On $[0,x]$, the functions $\frac{2t}{t^2-n^2}$ are continuous and bounded by $\frac2{n^2-1}$ for $n\ge2$, so the series
of Theorem~\ref{thm:cot} minus its first term converges uniformly on $[0,x]$, and the $n=1$ term is continuous there too.
Integrating termwise from $0$ to $x$ (Lemma~\ref{lem:uniform}(c)),
\[
\begin{aligned}\int_0^x\Bigl(\pi\cot\pi t-\frac1t\Bigr)\dd t&=\sum_{n\ge1}\int_0^x\frac{2t}{t^2-n^2}\dd t\\&=\sum_{n\ge1}\ln\Bigl(1-\frac{x^2}{n^2}\Bigr).\end{aligned}
\]
The integrand on the left is the derivative of $\ln\frac{\sin\pi t}{\pi t}$, which tends to $0$ at $0$, so the left side is
$\ln\frac{\sin\pi x}{\pi x}$. Exponentiating gives $\sin\pi x=P(x)$ for $0<x<1$; both sides are odd and vanish at $0$, so the
identity holds on $\left]-1,1\right[$.

To extend it, write the partial products as
\[
\begin{aligned}P_N(x)&=\pi x\prod_{n=1}^N\frac{(n-x)(n+x)}{n^2}\eqbr=\frac{(-1)^N\pi}{(N!)^2}\prod_{k=-N}^N(x+k).\end{aligned}
\]
Then, for $x\ne N$,
\[
\begin{aligned}P_N(x+1)&=\frac{(-1)^N\pi}{(N!)^2}\prod_{k=-N+1}^{N+1}(x+k)\eqbr=P_N(x)\,\frac{x+N+1}{x-N},\end{aligned}
\]
and letting $N\to\infty$, $P(x+1)=-P(x)$. Since $\sin\pi(x+1)=-\sin\pi x$ too, the identity on $\left]-1,1\right[$ propagates to
$\left]0,2\right[$, then to every interval $\left]k-1,k+1\right[$, hence to $\R$.
\end{proof}

\begin{figure}[ht]
\centering
\fig{ch03-sineproduct}
\caption{$\sin\pi x$ (thick) and the partial products $\pi x\prod_{n\le N}(1-\frac{x^2}{n^2})$ for $N=1,2,5$. Each partial product has
the right roots up to $\pm N$ and then escapes; the convergence is slow, with error of order $\frac{x^2}N$.}
\label{fig:sineproduct}
\end{figure}

\begin{corollary}[Wallis, 1656; Vi\`ete, 1593]\label{cor:wallis}
\[
\begin{aligned}\frac\pi2&=\prod_{n\ge1}\frac{(2n)^2}{(2n-1)(2n+1)}\eqbr=\frac21\cdot\frac23\cdot\frac43\cdot\frac45\cdot\frac65\cdot\frac67\cdots,\end{aligned}
\]
\[
\begin{aligned}\frac2\pi&=\prod_{k\ge1}\cos\frac\pi{2^{k+1}}\eqbr=\frac{\sqrt2}2\cdot\frac{\sqrt{2+\sqrt2}}2\cdot\frac{\sqrt{2+\sqrt{2+\sqrt2}}}2\cdots\end{aligned}
\]
\end{corollary}
\begin{proof}
At $x=\frac12$ the sine product gives $1=\frac\pi2\prod_n\bigl(1-\frac1{4n^2}\bigr)=\frac\pi2\prod_n\frac{(2n-1)(2n+1)}{(2n)^2}$, and the
product of the inverses of positive convergent factors is the inverse of the product. For Vi\`ete, iterate
$\sin y=2\sin\frac y2\cos\frac y2$: for $y\ne0$,
\[
\begin{gathered}\sin y=2^m\sin\frac y{2^m}\prod_{k=1}^m\cos\frac y{2^k},\\\prod_{k=1}^m\cos\frac y{2^k}=\frac{\sin y}{y}\cdot\frac{y/2^m}{\sin(y/2^m)}\to\frac{\sin y}y.\end{gathered}
\]
At $y=\frac\pi2$ this is $\frac2\pi$. Finally $\cos\frac\pi4=\frac{\sqrt2}2$ and $\cos\frac\theta2=\sqrt{\frac{1+\cos\theta}2}$, which gives the nested
radicals.
\end{proof}

Vi\`ete's formula, the first infinite product in the history of mathematics, has a geometric meaning
(Figure~\ref{fig:viete}): the area of the regular $2^{k+1}$-gon inscribed in the unit circle is
$2^{k}\sin\frac{2\pi}{2^{k+1}}$, and passing from $2^k$ to $2^{k+1}$ sides multiplies the area by $\frac1{\cos(\pi/2^k)}$. Starting
from the square, of area $2$, and approaching the disc, of area $\pi$, the product of the factors $\cos\frac\pi{2^k}$ ($k\ge2$)
is $\frac2\pi$.

\begin{figure}[ht]
\centering
\fig{ch03-viete}
\caption{Vi\`ete's product: inscribed regular polygons with $4$, $8$ and $16$ sides, of areas $2$, $2\sqrt2\approx2.83$ and
$3.06$; each doubling multiplies the area by $\frac1{\cos(\pi/2^k)}$, and the limit is $\pi$.}
\label{fig:viete}
\end{figure}

\section{Euler's formula for \texorpdfstring{$\zeta(2k)$}{zeta(2k)}}

Write $\zeta(s)=\sum_{n\ge1}n^{-s}$ for $s>1$.

\begin{theorem}[Euler, 1735--1740]\label{thm:zeta2k}
For every $k\ge1$,
\[
\zeta(2k)=\frac{(-1)^{k+1}B_{2k}(2\pi)^{2k}}{2\,(2k)!}.
\]
In particular $\zeta(2)=\frac{\pi^2}6$, $\zeta(4)=\frac{\pi^4}{90}$, $\zeta(6)=\frac{\pi^6}{945}$, $\zeta(8)=\frac{\pi^8}{9450}$.
\end{theorem}
\begin{proof}
\emph{The analytic side.} For $0<|x|<1$, Theorem~\ref{thm:cot} gives
\[
\begin{aligned}\pi x\cot\pi x&=1-2\sum_{n\ge1}\frac{x^2/n^2}{1-x^2/n^2}\eqbr=1-2\sum_{n\ge1}\sum_{k\ge1}\frac{x^{2k}}{n^{2k}}.\end{aligned}
\]
The double series converges absolutely, since $\sum_n\frac{x^2/n^2}{1-x^2/n^2}<\infty$, so the summations can be exchanged
(Lemma~\ref{lem:packets}): $\pi x\cot\pi x=Z(x)$, where $Z(x)=1-2\sum_{k\ge1}\zeta(2k)x^{2k}$ is a power series converging for
$|x|<1$. Multiply by $\frac{\sin\pi x}{\pi x}=S(x)=\sum_j\frac{(-1)^j(\pi x)^{2j}}{(2j+1)!}$, which converges everywhere: by the
Cauchy product of absolutely convergent series, $\cos\pi x=S(x)Z(x)$ as power series on $\left]-1,1\right[$, hence, by the
uniqueness of power series coefficients, as formal power series.

\emph{The formal side.} The proof of Proposition~\ref{prop:bernoulli}(b) showed the formal identity
$(e^t-1)f(t)=\frac t2(e^t+1)$, where $f(t)=\sum_kB_{2k}\frac{t^{2k}}{(2k)!}$. Substitute $t=2iy$ (formal series with complex
coefficients) and multiply by $e^{-iy}$, using $e^{a}e^{b}=e^{a+b}$:
\[
(e^{iy}-e^{-iy})f(2iy)=iy\,(e^{iy}+e^{-iy}).
\]
Comparing coefficients, $e^{iy}-e^{-iy}=2i\sin y$ and $e^{iy}+e^{-iy}=2\cos y$ as formal series, so
$\sin y\cdot f(2iy)=y\cos y$. With $y=\pi x$ and dividing by $\pi x$: $S(x)W(x)=\cos\pi x$, where
$W(x)=f(2\pi ix)=\sum_k(-1)^kB_{2k}\frac{(2\pi x)^{2k}}{(2k)!}$.

So $S\,Z=S\,W$ in $\C[[x]]$, and since $S$ has constant term $1$, it is invertible and $Z=W$. The coefficients of $x^{2k}$ give
$-2\zeta(2k)=(-1)^kB_{2k}\frac{(2\pi)^{2k}}{(2k)!}$. The values follow from $B_2=\frac16$, $B_4=-\frac1{30}$, $B_6=\frac1{42}$, $B_8=-\frac1{30}$.
\end{proof}

Two consequences are worth stating. Since $\zeta(2k)>0$, the signs of the $B_{2k}$ alternate: $(-1)^{k+1}B_{2k}>0$. And since
$1<\zeta(2k)\le\zeta(2)<2$,
\[
|B_{2k}|=\frac{2(2k)!}{(2\pi)^{2k}}\zeta(2k)\sim\frac{2(2k)!}{(2\pi)^{2k}},
\]
which grows faster than any exponential: this is the growth announced in Chapter~1. As for $\zeta(3)$, nobody knows a closed
form. Ap\'ery proved in 1978 that it is irrational, and it is unknown whether $\zeta(5)$ is.

\begin{remark}
The proof shows that $\pi x\cot\pi x$ is a power series on $\left]-1,1\right[$ and not beyond: its coefficients $-2\zeta(2k)$ do not
tend to $0$. In the variable $t=2\pi ix$, the series $\frac t{e^t-1}$ therefore has radius of convergence exactly $2\pi$, the
distance from $0$ to its nearest poles $\pm2\pi i$. Chapter~11 explains why the radius of a power series is always the distance
to the nearest singularity in the complex plane.
\end{remark}

\section{Irrationality of \texorpdfstring{$\pi$}{pi}}

Lambert proved in 1761 that $\pi$ is irrational, using a continued fraction for $\tan x$. The following short proof, in the
style of Hermite and Niven, gives more: $\pi^2$ is irrational, so $\pi$ is not even the square root of a rational.

\begin{theorem}\label{thm:piirrational}
$\pi^2$ is irrational. In particular $\pi$ is irrational.
\end{theorem}
\begin{proof}
For $n\ge1$ let $f(x)=\frac{x^n(1-x)^n}{n!}$. Then:
(i) $0<f(x)<\frac1{n!}$ for $0<x<1$;
(ii) every derivative $f^{(k)}(0)$ is an integer: writing $f=\frac1{n!}\sum_{i=n}^{2n}c_ix^i$ with integers $c_i$,
$f^{(k)}(0)=\frac{k!}{n!}c_k$, which is $0$ for $k<n$ and an integer for $k\ge n$;
(iii) every $f^{(k)}(1)$ is an integer, because $f(1-x)=f(x)$ gives $f^{(k)}(1)=(-1)^kf^{(k)}(0)$.

Suppose $\pi^2=\frac ab$ with positive integers $a$, $b$. Let
\[
F(x)=b^n\sum_{k=0}^n(-1)^k\pi^{2n-2k}f^{(2k)}(x).
\]
Since $b^n\pi^{2n-2k}=a^{n-k}b^k$ is an integer, (ii) and (iii) show that $F(0)$ and $F(1)$ are integers. Since $f^{(2n+2)}=0$
($f$ has degree $2n$), the sum telescopes in $F''+\pi^2F=b^n\pi^{2n+2}f=\pi^2a^nf$. Therefore the function $\Phi(x)=F'(x)\sin\pi x-\pi F(x)\cos\pi x$ satisfies
\[
\begin{aligned}\Phi'(x)&=\bigl(F''(x)+\pi^2F(x)\bigr)\sin\pi x\eqbr=\pi^2a^nf(x)\sin\pi x,\end{aligned}
\]
and integrating over $[0,1]$,
\[
\begin{aligned}\pi a^n\int_0^1f(x)\sin\pi x\dd x&=\Bigl[\frac{\Phi(x)}\pi\Bigr]_0^1\eqbr=F(1)+F(0)\in\Z,\end{aligned}
\]
since $\Phi(1)=\pi F(1)$ and $\Phi(0)=-\pi F(0)$. By (i), the left side lies strictly between $0$ and $\frac{\pi a^n}{n!}$. Since $\frac{\pi a^n}{n!}\to0$, choosing $n$ large makes it an
integer strictly between $0$ and $1$, a contradiction.
\end{proof}

The proof is a template: an integral that is a non-zero integer by an algebraic argument and smaller than $1$ by an analytic
one. Hermite used it in 1873 to prove that $e$ is transcendental, and Lindemann in 1882 to prove that $\pi$ is, which
settled the squaring of the circle.

\section{Exercises}

\begin{exercise}[$\star$]\label{ex:leibniz}
Show that $\frac1{1+x^2}=\sum_{k=0}^{n-1}(-1)^kx^{2k}+\frac{(-1)^nx^{2n}}{1+x^2}$, and deduce Leibniz's series
$\frac\pi4=1-\frac13+\frac15-\frac17+\cdots$, with an error at most $\frac1{2n+1}$ after $n$ terms.
\end{exercise}

\begin{exercise}[$\star$]\label{ex:cotsums}
Deduce from Theorem~\ref{thm:cot} that for real $a\notin\Z$,
\[
\sum_{n\ge1}\frac1{n^2-a^2}=\frac1{2a^2}-\frac{\pi\cot\pi a}{2a},
\]
and evaluate it at $a=\frac12$; check the result by telescoping. What happens as $a\to0$?
\end{exercise}

\begin{exercise}[$\star\star$]\label{ex:zetavalues}
Compute $\zeta(6)$ and $\zeta(8)$ from Theorem~\ref{thm:zeta2k}, and show that $\sum_{n\ge0}\frac1{(2n+1)^4}=\frac{\pi^4}{96}$ and
$\sum_{n\ge1}\frac{(-1)^{n-1}}{n^2}=\frac{\pi^2}{12}$.
\end{exercise}

\begin{exercise}[$\star\star$]\label{ex:ritelli}
\emph{Basel by a double integral.} Using Tonelli's theorem for non-negative continuous functions on $\left]0,\infty\right[^2$
(Admitted (proved in Volume~7)), compute $\int_0^\infty\!\int_0^\infty\frac{\mathrm dx\,\mathrm dy}{(1+y)(1+x^2y)}$ in both orders, and deduce
$\int_0^1\frac{-\ln x}{1-x^2}\dd x=\frac{\pi^2}8$, then $\zeta(2)=\frac{\pi^2}6$.
\end{exercise}

\begin{exercise}[$\star\star$]\label{ex:eirrational}
(a) Show that $e$ is irrational: if $e=\frac pq$, look at $q!\,e-\sum_{k\le q}\frac{q!}{k!}$.
(b) Show that $e$ is not a root of a quadratic $ae^2+be+c=0$ with integers $a,b,c$ not all zero. (Write $ae+b+\frac ce=0$
and multiply by $n!$ for a suitable $n$.)
\end{exercise}

\begin{exercise}[$\star\star\star$]\label{ex:sqrtperiod}
Let $d\ge2$ be an integer that is not a square, and $a_0=\lfloor\sqrt d\rfloor$. Volume~2, Chapter~8, showed that the continued fraction
of $\sqrt d$ is eventually periodic. Call a quadratic irrational $x$ \emph{reduced} if $x>1$ and its conjugate $x'$ satisfies
$-1<x'<0$.
\begin{enumerate}[label=(\alph*)]
\item Show that if $x$ is reduced, so are all its complete quotients $x_n$ (with $x_0=x$, $x_{n+1}=\frac1{x_n-\lfloor x_n\rfloor}$).
\item Show that for reduced $x_n$, $\lfloor x_n\rfloor=\bigl\lfloor-\frac1{x_{n+1}'}\bigr\rfloor$, so $x_n$ is determined by $x_{n+1}$. Deduce
that the continued fraction of a reduced quadratic irrational is purely periodic.
\item Apply this to $\sqrt d+a_0$ and deduce that $\sqrt d=[a_0;\overline{a_1,\dots,a_{\ell-1},2a_0}]$: the period ends with
$2a_0$.
\end{enumerate}
\end{exercise}

\begin{exercise}[$\star\star$]\label{pb:coprime52}
Show that
\[
\sum_{\substack{a,b\ge1\\\gcd(a,b)=1}}\frac1{a^2b^2}=\frac52.
\]
More generally, compute $\sum_{\gcd(a,b)=1}(ab)^{-s}$ for $s>1$, and interpret the result as a Dirichlet series
$\sum_nc_nn^{-s}$ with an explicit $c_n$.
\end{exercise}

\begin{exercise}[$\star\star$]\label{pb:levy}
\emph{How fast the denominators grow.} Let $x$ be uniform in $]0,1[$ and $\frac{p_n}{q_n}$ its convergents.
\begin{enumerate}[label=(\alph*)]
\item Show that $|q_nx-p_n|=x\,Tx\cdots T^nx$, and that this number lies between $\frac1{2q_{n+1}}$ and $\frac1{q_{n+1}}$.
\item Show that $\int_0^1\frac{-\log t}{1+t}\,dt=\frac{\pi^2}{12}$.
\item Using Problem~\ref{pb:gausskuzmin}, prove that
\[
\mathbb E[\log q_n]=\frac{\pi^2}{12\log2}\,n+O(1),\qquad\frac{\pi^2}{12\log2}=1.18656\ldots
\]
\end{enumerate}
\end{exercise}

\section{Hints}

\begin{hint}{ex:leibniz}
Integrate from $0$ to $1$; the remainder is at most $\int_0^1x^{2n}\dd x$.
\end{hint}
\begin{hint}{ex:cotsums}
$\frac{2a}{a^2-n^2}=-\frac{2a}{n^2-a^2}$. At $a=\frac12$, $\frac1{n^2-1/4}=\frac2{2n-1}-\frac2{2n+1}$. For $a\to0$, use
$\pi a\cot\pi a=1-\frac{\pi^2a^2}3+O(a^4)$.
\end{hint}
\begin{hint}{ex:zetavalues}
Remove the even terms: $\sum\frac1{(2n)^4}=\frac{\zeta(4)}{16}$. For the alternating sum, subtract twice the even terms.
\end{hint}
\begin{hint}{ex:ritelli}
In one order, $\int_0^\infty\frac{\mathrm dx}{1+x^2y}=\frac\pi{2\sqrt y}$. In the other, split
$\frac1{(1+y)(1+x^2y)}$ into partial fractions in $y$. Then send $\left]1,\infty\right[$ to $\left]0,1\right[$ by $x\mapsto\frac1x$, and expand
$\frac1{1-x^2}$.
\end{hint}
\begin{hint}{ex:eirrational}
In (a) the difference is an integer in $\left]0,1\right[$. In (b), $n!\,e$ and $\frac{n!}e$ are close to integers, with errors of
opposite signs or small enough; take $n$ large and look at the sizes of the remainders.
\end{hint}
\begin{hint}{ex:sqrtperiod}
(a) $x_{n+1}'=\frac1{x_n'-a_n}$. (b) From $x_n'=a_n+\frac1{x_{n+1}'}$, $-\frac1{x_{n+1}'}=a_n-x_n'\in\left]a_n,a_n+1\right[$. If
$x_m=x_{m+\ell}$, go backwards. (c) $\sqrt d+a_0$ has integer part $2a_0$ and the same fractional part as $\sqrt d$.
\end{hint}

\begin{hint}{pb:coprime52}
Every pair $(a,b)$ is $(ga',gb')$ with $g=\gcd(a,b)$ and $\gcd(a',b')=1$.
\end{hint}

\begin{hint}{pb:levy}
(a) $x=\frac{p_n+p_{n-1}T^nx}{q_n+q_{n-1}T^nx}$ (Volume~2, Proposition~8.1). (c) Take expectations in $\log q_{n+1}=-\sum_{j\le n}\log T^jx+O(1)$.
\end{hint}

\section{Solutions}

\begin{solution}{ex:leibniz}
\idea A geometric sum with an explicit remainder, integrated.
\writeup The identity is the geometric sum $\sum_{k<n}(-x^2)^k=\frac{1-(-x^2)^n}{1+x^2}$. Integrating over $[0,1]$,
\[
\begin{aligned}\frac\pi4&=\int_0^1\frac{\mathrm dx}{1+x^2}\\&=\sum_{k=0}^{n-1}\frac{(-1)^k}{2k+1}+(-1)^n\int_0^1\frac{x^{2n}}{1+x^2}\dd x,\end{aligned}
\]
and the remainder is at most $\int_0^1x^{2n}\dd x=\frac1{2n+1}$ in absolute value. It tends to $0$, which gives the series. The
error bound shows how slow it is: a million terms give about six digits.
\end{solution}

\begin{solution}{ex:cotsums}
\idea Theorem~\ref{thm:cot} is already this sum, with a sign and a factor.
\writeup For $a\notin\Z$, $\sum_{n\ge1}\frac{2a}{a^2-n^2}=\pi\cot\pi a-\frac1a$, so
$\sum_{n\ge1}\frac1{n^2-a^2}=-\frac1{2a}\bigl(\pi\cot\pi a-\frac1a\bigr)=\frac1{2a^2}-\frac{\pi\cot\pi a}{2a}$. At $a=\frac12$, $\cot\frac\pi2=0$, so the sum is
$2$; indeed $\sum_{n\ge1}\frac1{n^2-1/4}=\sum_{n\ge1}\bigl(\frac2{2n-1}-\frac2{2n+1}\bigr)=2$ by telescoping. As $a\to0$,
$\frac1{2a^2}-\frac{\pi\cot\pi a}{2a}=\frac{1-\pi a\cot\pi a}{2a^2}=\frac{\pi^2a^2/3+O(a^4)}{2a^2}\to\frac{\pi^2}6$, and the left side tends to
$\zeta(2)$ by uniform convergence near $0$: a second proof of Basel, from the cotangent alone.
\end{solution}

\begin{solution}{ex:zetavalues}
\idea Plug in the Bernoulli numbers; then separate odd and even $n$.
\writeup With $B_6=\frac1{42}$: $\zeta(6)=\frac{(2\pi)^6}{42\cdot2\cdot720}=\frac{64\pi^6}{60480}=\frac{\pi^6}{945}$. With $B_8=-\frac1{30}$:
$\zeta(8)=\frac{(2\pi)^8}{30\cdot2\cdot40320}=\frac{256\pi^8}{2419200}=\frac{\pi^8}{9450}$. The even terms of $\zeta(4)$ sum to $\frac{\zeta(4)}{16}$, so
$\sum\frac1{(2n+1)^4}=\frac{15}{16}\cdot\frac{\pi^4}{90}=\frac{\pi^4}{96}$. Finally $\sum\frac{(-1)^{n-1}}{n^2}=\zeta(2)-2\cdot\frac{\zeta(2)}4=\frac{\zeta(2)}2=\frac{\pi^2}{12}$;
the rearrangement is legitimate because the series converge absolutely.
\end{solution}

\begin{solution}{ex:ritelli}
\idea Compute one double integral twice. One order gives $\pi\cdot\frac\pi2$ through two arctangents; the other produces a
logarithm whose integral is a sum over odd squares.
\writeup Let $I$ be the integral of the positive continuous function $\frac1{(1+y)(1+x^2y)}$. First integrating in $x$: with
$x=\frac u{\sqrt y}$, $\int_0^\infty\frac{\mathrm dx}{1+x^2y}=\frac1{\sqrt y}\int_0^\infty\frac{\mathrm du}{1+u^2}=\frac\pi{2\sqrt y}$. Then with $y=t^2$,
$I=\frac\pi2\int_0^\infty\frac{\mathrm dy}{\sqrt y(1+y)}=\frac\pi2\int_0^\infty\frac{2\,\mathrm dt}{1+t^2}=\frac{\pi^2}2$. Now integrate in $y$ first, for $x\ne1$:
\[
\frac1{(1+y)(1+x^2y)}=\frac1{1-x^2}\Bigl(\frac1{1+y}-\frac{x^2}{1+x^2y}\Bigr),
\]
so $\int_0^\infty\frac{\mathrm dy}{(1+y)(1+x^2y)}=\frac1{1-x^2}\Bigl[\ln\frac{1+y}{1+x^2y}\Bigr]_0^\infty=\frac{\ln(1/x^2)}{1-x^2}=\frac{-2\ln x}{1-x^2}$, a
positive function (also at $x=1$ by continuity, where it is $1$). By Tonelli, $\int_0^\infty\frac{-2\ln x}{1-x^2}\dd x=\frac{\pi^2}2$. The
substitution $x=\frac1u$ maps $\int_1^\infty\frac{-\ln x}{1-x^2}\dd x$ to $\int_0^1\frac{\ln u}{1-u^{-2}}\cdot\frac{\mathrm du}{u^2}=\int_0^1\frac{-\ln u}{1-u^2}\dd u$, so
$\int_0^1\frac{-\ln x}{1-x^2}\dd x=\frac{\pi^2}8$. On $\left]0,1\right[$, $\frac{-\ln x}{1-x^2}=\sum_{n\ge0}(-\ln x)x^{2n}$ with non-negative terms, and
$\int_0^1(-\ln x)x^{2n}\dd x=\frac1{(2n+1)^2}$ (integrate by parts), so by monotone convergence (Volume~4, Chapter~5)
$\sum_{n\ge0}\frac1{(2n+1)^2}=\frac{\pi^2}8$. Since the even terms of $\zeta(2)$ sum to $\frac{\zeta(2)}4$, $\zeta(2)=\frac43\cdot\frac{\pi^2}8=\frac{\pi^2}6$.
\end{solution}

\begin{solution}{ex:eirrational}
\idea Multiplying the series of $e$ by $n!$ leaves an integer plus a tail of size about $\frac1n$. For (b), one needs the tails
of $e$ and $\frac1e$ at the same time, and the alternating tail of $\frac1e$ has a known sign.
\writeup (a) If $e=\frac pq$, then $N=q!\,e-\sum_{k=0}^q\frac{q!}{k!}$ is an integer. But
$N=\sum_{k>q}\frac{q!}{k!}=\frac1{q+1}+\frac1{(q+1)(q+2)}+\cdots<\sum_{j\ge1}\frac1{(q+1)^j}=\frac1q\le1$ and $N>0$, a contradiction.

(b) Suppose $ae^2+be+c=0$, so $ae+b+\frac ce=0$, with $a$, $c$ not both zero (if $a=c=0$ then $b=0$). Replacing the triple by its
opposite, we may assume $a\ge0$; if $a=0$ then $c\ne0$ and $\frac ce=-b$ contradicts the irrationality of $e$ (a). So $a>0$. For
$n\ge1$ write $n!\,e=A_n+r_n$ and $n!\,e^{-1}=C_n+s_n$ with integers $A_n=\sum_{k\le n}\frac{n!}{k!}$,
$C_n=\sum_{k\le n}\frac{(-1)^kn!}{k!}$ and $0<r_n<\frac1n$ by (a); for the alternating tail
$s_n=(-1)^{n+1}\bigl(\frac1{n+1}-\frac1{(n+1)(n+2)}+\cdots\bigr)$, so $s_n$ has the sign of $(-1)^{n+1}$ and $|s_n|<\frac1{n+1}$. Multiply the
relation by $n!$: $aA_n+bn!+cC_n=-(ar_n+cs_n)$, an integer. If $c\ge0$, choose $n$ odd, so $s_n>0$; if $c<0$, choose $n$ even, so
$s_n<0$ and $cs_n>0$. Either way $ar_n+cs_n>0$, and it is less than $\frac{a+|c|}n<1$ for $n$ large. So a positive integer is less
than $1$, a contradiction.
\end{solution}

\begin{solution}{ex:sqrtperiod}
\idea Reduced numbers are closed under the continued fraction step, and the step can be inverted by looking at conjugates;
a map that is eventually periodic and invertible is periodic from the start.
\writeup (a) Let $x_n$ be reduced and $a_n=\lfloor x_n\rfloor\ge1$. Then $x_{n+1}=\frac1{x_n-a_n}>1$, and conjugating (the conjugation is a field
automorphism of $\Q(\sqrt d)$ fixing $\Q$), $x_{n+1}'=\frac1{x_n'-a_n}$. Since $x_n'-a_n<-1$, $-1<x_{n+1}'<0$.

(b) From $x_n=a_n+\frac1{x_{n+1}}$, conjugation gives $x_n'=a_n+\frac1{x_{n+1}'}$, so $-\frac1{x_{n+1}'}=a_n-x_n'$, which lies in $\left]a_n,a_n+1\right[$
since $-1<x_n'<0$. Hence $a_n=\bigl\lfloor-\frac1{x_{n+1}'}\bigr\rfloor$, and $x_n=a_n+\frac1{x_{n+1}}$ is a function of $x_{n+1}$. By Volume~2, the
sequence $(x_n)$ is eventually periodic: $x_m=x_{m+\ell}$ for some $m\ge0$, $\ell\ge1$. If $m\ge1$, applying the inverse step gives
$x_{m-1}=x_{m+\ell-1}$; repeating, $x_0=x_\ell$, so the expansion is purely periodic.

(c) Let $x=\sqrt d+a_0$. Then $x>1$ and $x'=a_0-\sqrt d\in\left]-1,0\right[$, so $x$ is reduced, and its expansion is purely periodic,
$x=[\overline{b_0;b_1,\dots,b_{\ell-1}}]$. Its integer part is $b_0=2a_0$, and $x-2a_0=\sqrt d-a_0$, so the complete quotients of $x$ from
the first on are those of $\sqrt d$: $b_k=a_k$ for $k\ge1$. Hence $\sqrt d=x-a_0=[a_0;a_1,\dots,a_{\ell-1},2a_0,a_1,\dots]$, that is,
$\sqrt d=[a_0;\overline{a_1,\dots,a_{\ell-1},2a_0}]$. For example $\sqrt{14}=[3;\overline{1,2,1,6}]$.
\end{solution}

\begin{solution}{pb:coprime52}
\idea Pull out the gcd: summing over all pairs is summing over $g$ and over coprime pairs, and the gcd contributes a factor
$\zeta(2s)$. So the coprime sum is $\frac{\zeta(s)^2}{\zeta(2s)}$, and at $s=2$ the values of $\zeta(2)$ and $\zeta(4)$ do the
rest.

\writeup Let $s>1$. All the series have positive terms, so they may be rearranged freely. Every pair $(a,b)$ of positive
integers is uniquely $(ga',gb')$ with $g=\gcd(a,b)$ and $\gcd(a',b')=1$. Hence
\[
\zeta(s)^2=\sum_{a,b\ge1}\frac1{(ab)^s}=\sum_{g\ge1}\frac1{g^{2s}}\sum_{\gcd(a',b')=1}\frac1{(a'b')^s}=\zeta(2s)\,S(s),
\]
so $S(s)=\frac{\zeta(s)^2}{\zeta(2s)}$. With $\zeta(2)=\frac{\pi^2}6$ and $\zeta(4)=\frac{\pi^4}{90}$ (Volume~5, Chapter~3),
$S(2)=\frac{\pi^4}{36}\cdot\frac{90}{\pi^4}=\frac52$.

Grouping by the product $n=a'b'$, $S(s)=\sum_nc_nn^{-s}$ where $c_n$ is the number of ordered coprime pairs $(a',b')$ with
$a'b'=n$. Coprimality forces each prime power $p^e\parallel n$ to go entirely to $a'$ or to $b'$, so $c_n=2^{\omega(n)}$, where
$\omega(n)$ is the number of distinct prime factors of $n$. So $\sum_n2^{\omega(n)}n^{-s}=\frac{\zeta(s)^2}{\zeta(2s)}$.
\tools Volume~2, Chapters~1 and~6 (gcd, multiplicative functions); Volume~5, Chapter~3 ($\zeta(2)$, $\zeta(4)$).
\end{solution}

\begin{solution}{pb:levy}
\idea The product $x\,Tx\cdots T^nx$ is both the approximation error of the $n$-th convergent and a product of the values of a
function along the orbit; taking logarithms turns it into a sum, whose terms have, by Gauss--Kuzmin, almost the same law.

\writeup (a) The complete quotient after $a_n$ is $1/T^nx$, so by Volume~2, Proposition~8.1, $x=\frac{p_n/T^nx+p_{n-1}}{q_n/T^nx+q_{n-1}}=\frac{p_n+p_{n-1}T^nx}{q_n+q_{n-1}T^nx}$.
Hence, with $p_nq_{n-1}-p_{n-1}q_n=\pm1$,
\[
\beta_n:=|q_nx-p_n|=\frac{T^nx}{q_n+q_{n-1}T^nx}.
\]
Applied at rank $n-1$, and using $T^{n-1}x=\frac1{a_n+T^nx}$ and $q_n=a_nq_{n-1}+q_{n-2}$, this gives $\beta_{n-1}=\frac1{q_n+q_{n-1}T^nx}$;
so $\beta_n=T^nx\cdot\beta_{n-1}$. With $\beta_0=x$ (as $p_0=0$, $q_0=1$), $\beta_n=x\,Tx\cdots T^nx$, and the formula for $\beta_n$ at rank
$n+1$, $\beta_n=\frac1{q_{n+1}+q_nT^{n+1}x}$, puts it between $\frac1{2q_{n+1}}$ and $\frac1{q_{n+1}}$ (as $q_n\le q_{n+1}$).

(b) $\int_0^1(-\log t)t^m\,dt=\frac1{(m+1)^2}$, and expanding $\frac1{1+t}=\sum_m(-t)^m$ (the partial sums are dominated by $1$),
$\int_0^1\frac{-\log t}{1+t}dt=\sum_{m\ge0}\frac{(-1)^m}{(m+1)^2}=\frac{\pi^2}{12}$.

(c) By (a), $\log q_{n+1}=\sum_{j=0}^n(-\log T^jx)+\theta$ with $-\log2\le\theta\le0$. By Problem~\ref{pb:gausskuzmin}(d), $T^jx$ has a density
$f_j$ with $|f_j(t)-\frac1{(1+t)\log2}|\le c^j$, so
\[
\Bigl|\mathbb E[-\log T^jx]-\frac{\pi^2}{12\log2}\Bigr|\le c^j\int_0^1(-\log t)\,dt=c^j,
\]
using (b). Summing over $j\le n$, $\mathbb E[\log q_{n+1}]=(n+1)\frac{\pi^2}{12\log2}+O\bigl(\frac1{1-c}+\log2\bigr)$.

\tools Volume~2, Chapter~8 (convergents); Volume~5, Chapter~3 (series for $\frac{\pi^2}{12}$); Problem~\ref{pb:gausskuzmin}.
\leads Lévy (1929) proved the almost sure version $\frac1n\log q_n\to\frac{\pi^2}{12\log2}$. For Euclid's algorithm on rationals,
Heilbronn (1969) and Dixon showed that the mean number of division steps for $\frac ab$ with $a<b$ is $\frac{12\log2}{\pi^2}\log b+O(1)$:
the reciprocal of Lévy's constant, as (c) suggests, since each step multiplies the denominator by about $e^{1.1866}$. Porter found
the next term.
\end{solution}

\section*{Sources and further reading}

\begin{itemize}
\item M.~Aigner and G.~Ziegler, \emph{Proofs from THE BOOK}, 6th ed., Springer, 2018: the chapters ``Three times $\pi^2/6$'',
``Cotangent and the Herglotz trick'' and ``Some irrational numbers''.
\item L.~Euler, \emph{De summis serierum reciprocarum}, Commentarii academiae scientiarum Petropolitanae 7 (1740), 123--134
(the Euler Archive, E41). W.~Dunham, \emph{Euler: The Master of Us All}, MAA, 1999, Chapter~3.
\item R.~Remmert, \emph{Theory of Complex Functions}, Springer, 1991, Chapter~11 (the cotangent and the sine product), and
\emph{Classical Topics in Complex Function Theory}, Springer, 1998, Chapter~1 (infinite products).
\item I.~Niven, \emph{A simple proof that $\pi$ is irrational}, Bulletin of the American Mathematical Society 53 (1947), 509.
\item D.~Ritelli, \emph{Another proof of $\zeta(2)=\frac{\pi^2}6$ using double integrals}, American Mathematical Monthly 120
(2013), 642--645.
\item R.~Ap\'ery, \emph{Irrationalit\'e de $\zeta(2)$ et $\zeta(3)$}, Ast\'erisque 61 (1979), 11--13; A.~van der Poorten, \emph{A proof that
Euler missed}, Mathematical Intelligencer 1 (1979), 195--203.
\item J.~Borwein and P.~Borwein, \emph{Pi and the AGM}, Wiley, 1987, Chapter~11 (irrationality and transcendence of $\pi$).
\item O.~Perron, \emph{Die Lehre von den Kettenbr\"uchen}, 1913, for Galois's theorem on purely periodic continued fractions.
\end{itemize}
